% Part: first-order-logic
% Chapter: completeness
% Section: lindenbaums-lemma

\documentclass[../../../include/open-logic-section]{subfiles}

\begin{document}

\iftag{FOL}
      {\olfileid{fol}{com}{lin}}
      {\olfileid{pl}{com}{lin}}

\olsection{લિન્ડનબાઉમનું સહાયક પ્રમેય}

\begin{explain}
હવે આપણે એવું સહાયક પ્રમેય સાબિત કરીશું જે બતાવે છે કે
!!{sentence}sનો કોઈ પણ સુસંગત ગણ એવા કોઈ વાક્યગણમાં સમાયેલો છે જે
માત્ર સુસંગત જ નહિ, પણ !!{complete} પણ છે. સાબિતીમાં એક સમયે એક
!!{sentence} ઉમેરીને દરેક પગલે ગણ સુસંગત રહે તેની ખાતરી કરવામાં આવે
છે. આ એવું કરીએ છીએ કે દરેક~$!A$ માટે કોઈક તબક્કે કાં તો~$!A$ અથવા
$\lnot !A$ ઉમેરાય. પછી આ રચનાના બધા તબક્કાઓના સંઘમાં કાં તો~$!A$ અથવા
તેનો નિષેધ~$\lnot !A$ હોય છે અને તેથી તે પૂર્ણ છે. તે સુસંગત પણ છે,
કારણ કે દરેક તબક્કે અસુસંગતતા દાખલ ન થાય તેની આપણે ખાતરી કરીએ છીએ.
\end{explain}

\begin{lem}[લિન્ડનબાઉમનું સહાયક પ્રમેય]
\ollabel{lem:lindenbaum} ભાષા~$\Lang{L}$ના દરેક સુસંગત ગણ~$\Gamma$ને
!!a{complete} અને સુસંગત ગણ~$\Gamma^*$ સુધી વિસ્તારી શકાય છે.
\end{lem}

\begin{proof}
ધારો કે~$\Gamma$ સુસંગત છે. ~ $\Lang L$નાં બધાં !!{sentence}sની
કોઈ પરિગણના $!A_0$, $!A_1$, \dots{} લો. $\Gamma_0=\Gamma$ વ્યાખ્યાયિત
કરીને મૂકો
\[
\Gamma_{n+1} =
\begin{cases}
\Gamma_n \cup \{ !A_n \} & \textrm{જો $\Gamma_n \cup \{!A_n\}$
  સુસંગત હોય;} \\
\Gamma_n \cup \{ \lnot !A_n \} & \textrm{અન્યથા.}
\end{cases}
\]
અને $\Gamma^*=\bigcup_{n\geq 0}\Gamma_n$ લો.

દરેક~$\Gamma_n$ સુસંગત છે: વ્યાખ્યા મુજબ~$\Gamma_0$ સુસંગત છે. જો
$\Gamma_{n+1}=\Gamma_n\cup\{!A_n\}$ હોય, તો તેનું કારણ પાછળનો ગણ
સુસંગત છે તે છે. જો તે સુસંગત ન હોય, તો
$\Gamma_{n+1}=\Gamma_n\cup\{\lnot !A_n\}$. આપણે ચકાસવું પડશે કે
$\Gamma_n\cup\{\lnot !A_n\}$ સુસંગત છે. ધારો કે તે સુસંગત નથી. તો
$\Gamma_n\cup\{!A_n\}$ અને $\Gamma_n\cup\{\lnot !A_n\}$ \emph{બંને}
અસુસંગત છે. આનો અર્થ નીચેના પરિણામ મુજબ~$\Gamma_n$ અસુસંગત થાય:
\iftag{FOL}{%
  \tagrefs{prfAX/{fol:axd:prv:prop:provability-exhaustive},
    prfSC/{fol:seq:prv:prop:provability-exhaustive},
    prfND/{fol:ntd:prv:prop:provability-exhaustive},
    prfTab/{fol:tab:prv:prop:provability-exhaustive}}}{%
  \tagrefs{prfAX/{pl:axd:prv:prop:provability-exhaustive},
    prfSC/{pl:seq:prv:prop:provability-exhaustive},
    prfND/{pl:ntd:prv:prop:provability-exhaustive},
    prfTab/{pl:tab:prv:prop:provability-exhaustive}}},
જે અનુમાનની પૂર્વધારણાનો વિરોધાભાસ છે.

દરેક~$n$ અને દરેક $i<n$ માટે $\Gamma_i\subseteq\Gamma_n$. આ~$n$ પરના
સરળ અનુમાનથી ફલિત થાય છે. $n=0$ માટે કોઈ $i<0$ નથી, તેથી દાવો આપોઆપ
સાચો છે. અનુમાનના પગલા માટે ધારો કે તે~$n$ માટે સાચો છે. જો $i<n+1$
હોય, તો $\Gamma_i\subseteq\Gamma_{n+1}$ એવું બતાવીએ. રચના મુજબ
$\Gamma_{n+1}=\Gamma_n\cup\{!A_n\}$ અથવા
$=\Gamma_n\cup\{\lnot !A_n\}$ છે. તેથી
$\Gamma_n\subseteq\Gamma_{n+1}$. જો $i<n+1$, તો અનુમાનની પૂર્વધારણા
મુજબ $i<n$ હોય ત્યારે $\Gamma_i\subseteq\Gamma_n$, અને $i=n$ હોય
ત્યારે સ્વાભાવિક રીતે $\Gamma_n\subseteq\Gamma_n$. હવે~$\subseteq$ની
પરંપરિતતા મુજબ $\Gamma_i\subseteq\Gamma_{n+1}$ મળે છે.

આના પરથી~$\Gamma^*$ સુસંગત છે તેમ ફલિત થાય છે. તેનું કારણ આ છે:
કોઈ સાન્ત $\Gamma'\subseteq\Gamma^*$ લો. જો~$\Gamma'$ ખાલી હોય, તો તે
સુસંગત છે. અન્યથા, દરેક $!B\in\Gamma'$ કોઈક~$i$ માટે~$\Gamma_i$માં
પણ હોય છે. આ સાન્ત સંખ્યામાં મળતા સૂચકોમાં સૌથી મોટો~$n$ લો. $i\le n$ હોય ત્યારે
$\Gamma_i\subseteq\Gamma_n$ હોવાથી દરેક $!B\in\Gamma'$ પણ
$\in\Gamma_n$, એટલે $\Gamma'\subseteq\Gamma_n$; અને~$\Gamma_n$
સુસંગત છે. આમ, દરેક સાન્ત ઉપગણ $\Gamma'\subseteq\Gamma^*$ સુસંગત છે. હવે
\iftag{FOL}{%
  \tagrefs{prfAX/{fol:axd:ptn:prop:proves-compact},
    prfSC/{fol:seq:ptn:prop:proves-compact},
    prfND/{fol:ntd:ptn:prop:proves-compact},
    prfTab/{fol:tab:ptn:prop:proves-compact}}}{%
  \tagrefs{prfAX/{pl:axd:ptn:prop:proves-compact},
    prfSC/{pl:seq:ptn:prop:proves-compact},
    prfND/{pl:ntd:ptn:prop:proves-compact},
    prfTab/{pl:tab:ptn:prop:proves-compact}}} મુજબ~$\Gamma^*$ સુસંગત
છે.

\sourcecorrection{OLCO-008}{મૂળની \texttt{lindenbaums-lemma.tex},
પંક્તિઓ~82--84માં મનસ્વી સાન્ત~$\Gamma'$માંથી મળતા સૂચકોમાં સૌથી મોટો
સૂચક લેવામાં આવે છે, પરંતુ ખાલી~$\Gamma'$ માટે એવો સૂચક હોતો નથી.
અહીં ખાલી કિસ્સો અલગથી નોંધ્યો છે.}

~$\Frm[L]$નું દરેક !!{sentence}~$\Gamma^*$ને વ્યાખ્યાયિત કરવા વપરાયેલી
યાદીમાં આવે છે. જો $!A_n\notin\Gamma^*$ હોય, તો તેનું કારણ
$\Gamma_n\cup\{!A_n\}$ અસુસંગત હતું એ છે. પરંતુ પછી $\lnot !A_n\in
\Gamma^*$; તેથી~$\Gamma^*$ !!{complete} છે.
\end{proof}

\end{document}
